\section*{Letter of Transmittal}
\addcontentsline{toc}{section}{Letter of Transmittal}
\thispagestyle{plain}

\begin{singlespace}
\noindent
California Polytechnic State University\\
Biological Sciences Department\\
1 Grand Avenue\\
San Luis Obispo, CA 93407-0830
\end{singlespace}

\vspace{2em}

\noindent December 31, 2025

\vspace{2em}

\begin{singlespace}
\noindent
Emma Pelton\\
The Xerces Society for Invertebrate Conservation\\
1631 NE Broadway, \#821\\
Portland, OR 97232-1425
\end{singlespace}

\vspace{2em}

\noindent\textbf{Re: Technical Report -- Subaward CalPol23-4132}

\vspace{1.5em}

\noindent Dear Ms.\ Pelton,

\vspace{1.5em}

\noindent This report constitutes the final technical deliverable under Subaward CalPol23-4132, ``Wind Model Validation for Monarch Butterfly Overwintering Groves,'' funded through USDA Forest Service Award 21-DG-11132762-274 via The Xerces Society for Invertebrate Conservation.

The original statement of work proposed computational fluid dynamics (CFD) simulations using LiDAR-derived forest geometry to model wind patterns in overwintering groves. During the first year of the project, the drone-mounted LiDAR platform required for aerial data acquisition was damaged during training operations and remained out of service for an extended period, making the originally proposed methodology infeasible within the project timeline. In coordination with the Xerces Society, the research approach was revised to pursue direct empirical investigation of wind effects on overwintering monarch butterfly clusters at Vandenberg Space Force Base.

The revised study deployed time-lapse cameras and wind sensors at monarch cluster roost sites over the 2023--2024 overwintering season, producing 1,894 paired observations across 78 days. This approach enabled the first direct empirical test of the wind disruption component of the microclimate hypothesis, which holds that wind speeds at or above 2 m/s physically dislodge butterflies from roosting clusters. This assumption has guided monarch habitat management for over three decades but had never been tested through direct observation of butterfly behavioral responses to wind.

The results do not support a wind disruption threshold. Monarch clusters persisted throughout the monitoring period despite routinely experiencing wind speeds well above 2 m/s, and statistical analysis found no meaningful relationship between wind speed and cluster dynamics. Instead, interactions between wind and sunlight emerged as significant predictors, suggesting that thermal regulation rather than wind avoidance may be the primary driver of butterfly movement at roost sites. These findings have direct implications for habitat management practices at overwintering sites managed by the Department of Defense, U.S.\ Forest Service, California State Parks, and other land managers, particularly regarding wind buffer zone requirements and habitat suitability criteria.

This research was conducted as a master's thesis at California Polytechnic State University and has undergone academic committee review. A peer-reviewed manuscript based on this work is being prepared for publication in \textit{Insects}. The enclosed report presents the complete methodology, results, and analysis in publication-quality format as specified in the statement of work. All data and analytical code produced under this subaward are available upon request.

Questions regarding this report may be directed to Kyle Nessen, M.S., Biological Sciences Department, California Polytechnic State University, San Luis Obispo, CA 93407.

\vspace{2em}

\begin{samepage}
\begin{singlespace}

\noindent Sincerely,

\vspace{3em}

\noindent\rule{0.45\textwidth}{0.4pt}\\[0.5em]
Francis Villablanca, Ph.D.\\
Principal Investigator\\
Professor Emeritus of Biology\\
California Polytechnic State University

\vspace{0.5em}

\noindent Date: \rule{0.3\textwidth}{0.4pt}

\vspace{3em}

\noindent\rule{0.45\textwidth}{0.4pt}\\[0.5em]
Kyle Nessen, M.S.\\
Graduate Researcher\\
Biological Sciences Department\\
California Polytechnic State University

\vspace{0.5em}

\noindent Date: \rule{0.3\textwidth}{0.4pt}

\end{singlespace}
\end{samepage}
